% =====================================================================
%  Kết luận (P1)
%  Đã viết: Phần B, C1–C4. Chờ số liệu cuối của P2 (A, C5) và P3 (D)
%  để điền các chỗ \todo.
% =====================================================================
\section{Kết luận}\label{sec:conclusion}

Báo cáo đã đi trọn chuỗi từ cơ sở lý thuyết của SVM và phương pháp hàm nhân, qua cài đặt từ đầu,
tới thực nghiệm và ứng dụng. Các kết quả chính theo từng phần như sau.

\paragraph{Lý thuyết (Phần B).} Mục~\ref{sec:theory} chứng minh đầy đủ sáu kết quả nền tảng. Bài toán
cực đại lề tương đương quy hoạch toàn phương lồi $\min \frac12\norm{\vw}^2$ và có nghiệm duy nhất; bài
toán đối ngẫu của SVM lề mềm chỉ phụ thuộc dữ liệu qua hàm nhân, và điều kiện KKT phân loại mỗi mẫu
thành ngoài lề, trên lề hoặc vi phạm lề; SVM lề mềm chính là cực tiểu hóa rủi ro hinge có chính quy
hóa với $C = 1/(\lambda n)$; định lý biểu diễn giải thích vì sao nghiệm chỉ cần $n$ hệ số; tổng, tích,
giới hạn và hàm mũ của hàm nhân xác định dương vẫn xác định dương, kéo theo tính hợp lệ của hàm nhân
Gauss; và tỉ lệ vector hỗ trợ chặn trên sai số bỏ-một. Trong quá trình chứng minh, hai giả thiết ngầm
của đề bài được nêu rõ: mỗi lớp có ít nhất một mẫu (B1, B6) và giá trị tối ưu hữu hạn (B4).

\paragraph{Cài đặt (Phần A).} \todo{P2: tóm tắt 2--3 câu kết quả của A1--A4: SMO, DCD, KRR, đa lớp,
Platt scaling; mức sai lệch so với LIBSVM (hàm mục tiêu đối ngẫu, tập vector hỗ trợ, dự đoán).}

\paragraph{Bài tập và thực nghiệm (Phần C).} Lời giải tay của C1 ($\vw = (1, 1)$, $b = -1$,
$\valpha = (\frac12, \frac12, 1)$, độ rộng lề $\sqrt2$) và C2 (siêu phẳng tách XOR trong không gian
đặc trưng của hàm nhân đa thức bậc hai, trùng với nghiệm lề cực đại $f(\vx) = -x_1x_2$) được SMO của
nhóm và LIBSVM tìm lại chính xác. Khảo sát $C$ và $\gamma$ (C4) bằng kiểm định chéo lồng nhau cho sai
số kiểm tra $7{,}0 \pm 1{,}1\%$ trên moons và $1{,}9 \pm 1{,}0\%$ trên breast cancer. Thực nghiệm xác
nhận các dự đoán của lý thuyết: khi $C$ tăng, độ rộng lề và số vector hỗ trợ đều không tăng với mọi
$\gamma$; $\gamma$ quá lớn dẫn tới ma trận Gram gần đơn vị và mô hình chỉ ghi nhớ dữ liệu huấn luyện;
và cận $\text{LOO} \le \nsv/n$ đúng ở cả $62$ cấu hình đã thử. \todo{P2: 1--2 câu kết quả C5
(thời gian SMO so với DCD khi $n$ tăng).}

\paragraph{Ứng dụng (Phần D).} \todo{P3: tóm tắt 2--3 câu: dữ liệu thu thập được, mô hình tốt nhất
và sai số khi chia theo phiên chụp, độ bền khi đổi thiết bị, khả năng dùng để kiểm tra nhãn gạo đặc sản.}

\paragraph{Hạn chế.}
\begin{itemize}
  \item Cận bỏ-một của B6 thường lỏng: trên các cấu hình đã thử, sai số bỏ-một thực tế chỉ bằng khoảng
        $0{,}36$--$0{,}37$ lần cận (trung vị). Với SVM lề mềm, cận chỉ được bảo đảm khi nghiệm có vector
        hỗ trợ tự do; khi mọi $\alpha_i = C$, sai số bỏ-một có thể lên tới $100\%$ dù sai số kiểm định
        chéo chỉ khoảng $17\%$.
  \item Khảo sát C4 dùng hai bộ dữ liệu chuẩn và lưới thô ($6 \times 5$ giá trị); chưa chạy trên đặc
        trưng hạt gạo của Phần D.
  \item \todo{P2, P3: hạn chế của phần mình (ví dụ: giới hạn thời gian ở C5, kích thước và độ đa dạng
        của dữ liệu gạo).}
\end{itemize}

\paragraph{Hướng phát triển.}
\begin{itemize}
  \item Thay cận $\nsv/n$ bằng các cận chặt hơn nhưng vẫn tính được sau một lần huấn luyện, như cận
        span hoặc cận bán kính--lề, để dùng trực tiếp cho việc chọn siêu tham số.
  \item Lặp lại khảo sát C4 trên đặc trưng hạt gạo để xem vùng $(C, \gamma)$ tốt có ổn định giữa các
        phiên chụp và thiết bị hay không.
  \item \todo{P2, P3: hướng phát triển cho phần cài đặt và ứng dụng.}
\end{itemize}
